% Section 4: Methodology

\section{Methodology}
\label{sec:methodology}

We follow MUSCAT's four phases \cite{muscat2023} with three changes: OSM building geometry, deterministic ray-traced line of sight, and separate cooperative and non-cooperative threat-weighted coverage. Phase~4 searches the resulting objective space with NSGA-II; the search procedure is a tool, and its fronts are checked against uniform random search at the same budget.

\subsection{Phase 1: Environment Setup}

Unlike MUSCAT's terrain models from DTED maps \cite{muscat2023}, we take building footprints from OpenStreetMap (OSM) \cite{osm_applications,elias2023exploring}: buildings within a radius of a given latitude and longitude are downloaded, heights are read from metadata or estimated, and geometries are reprojected to local UTM coordinates. The scene is then discretized into a 3D occupancy grid, \(\text{Grid}(i,j,k)=1\) if the voxel is occupied by a building and 0 otherwise, at resolution \(r\) (30--40\,m in this study).

\subsection{Phase 2: Sensor Attributes and ROI Definition}

\subsubsection{Sensor Models with Field-of-View Constraints}

Each sensor type has a detection model gated by an elevation window \(\theta_{elev}\in[\theta_{min},\theta_{max}]\), where \(\theta_{elev}\) is the elevation angle from sensor to target:

\textbf{Radar}: 360\textdegree{} horizontal coverage, $\theta_{elev} \in [-5^{\circ}, 30^{\circ}]$; \(P_D\) from the radar equation in transmit power, gain, target RCS, range, and wavelength.

\textbf{RF passive (RID equivalent)}: nearly omnidirectional, $\theta_{elev} \in [-30^{\circ}, 70^{\circ}]$, sensitivity \(-90\)\,dBm at 900\,MHz; \(P_D\) from received power against sensitivity.

\textbf{Electro-Optical (EO)}: 360\textdegree{} horizontal coverage, $\theta_{elev} \in [-20^{\circ}, 45^{\circ}]$, maximum range 150\,m.

\textbf{Acoustic}: 360\textdegree{} horizontal coverage, $\theta_{elev} \in [-20^{\circ}, 70^{\circ}]$, source level 80\,dB SPL, SNR threshold 6\,dB, maximum range 300\,m.

These values are those of the experiment configurations (Section~\ref{sec:setup}). Elevation limits prevent rooftop sensors from detecting targets far below the horizon. Every modality, including Acoustic, is gated by the same binary LoS test; this is conservative for acoustic sensing, which can partly propagate around obstacles.

Candidate sites for real scenes are generated on a grid with spacing \(s\), keeping points outside building footprints and assigning mounting heights \(h\sim U(25,35)\)\,m.

\subsection{Phase 3: Coverage Computation}

\subsubsection{Deterministic Ray-Tracing}

We replace statistical LoS probability models \cite{3gpp38901,itur_p1411} with deterministic ray-tracing on the occupancy grid. For a sensor at \(\mathbf{s}\) and a target voxel at \(\mathbf{t}\), \(P_D=0\) if the elevation angle falls outside the sensor's window; otherwise the segment from \(\mathbf{s}\) to \(\mathbf{t}\) is sampled voxel by voxel, and \(P_{LoS}=0\) if any sampled voxel is occupied and \(P_{LoS}=1\) otherwise. Coverage is evaluated over the whole free volume and, separately, on three airways: 20\,m (between buildings), 45\,m (above low-rise buildings), and 65\,m (above most buildings).

\subsubsection{Network Detection Probability}

For multiple sensors, we employ soft voting fusion:

\begin{equation}
\pnet = 1 - \prod_{i=1}^{N_s} (1 - P_{D,i})
\label{eq:pnet}
\end{equation}

where $P_{D,i}$ is the detection probability of the $i$-th sensor.

\subsubsection{Dual-Layer Threat-Weighted Coverage}
\label{sec:dual_metrics}

Critical assets \(a\) with location \(\mathbf{p}_a\), protection radius \(\rho_a\), and weight multiplier \(m_a\) define a threat map over free voxels \(v\):
\begin{equation}
w(v) = \max_a\; 100\, m_a \exp\!\left(-\frac{\lVert \mathbf{x}_v - \mathbf{p}_a \rVert^2}{2\rho_a^2}\right),
\label{eq:threat}
\end{equation}
with \(w(v)=0\) inside buildings. The network probability \(\pnet\) is computed twice: over all active sensors (cooperative layer, \(\pnet^{\mathrm{coop}}\)) and over Radar, EO, and Acoustic only (non-cooperative layer, \(\pnet^{\mathrm{nc}}\)), because RF receivers detect only targets that broadcast. Threat-weighted coverage is
\begin{equation}
M_{wp,\ell} = \frac{\sum_v w(v)\,\pnet^{\ell}(v)}{\sum_v w(v)}, \qquad \ell\in\{\mathrm{coop},\mathrm{nc}\}.
\label{eq:mwp}
\end{equation}
For airway analysis (Fig.~\ref{fig:airways}), the sums are restricted to the voxel layer containing the airway altitude.

\subsection{Phase 4: Search with NSGA-II}

The objectives are \(M_{wp,\mathrm{coop}}\), \(M_{wp,\mathrm{noncoop}}\), and CapEx (Eqs.~\eqref{eq:fitness}--\eqref{eq:mwp}), with no scalarization weights. We search them with NSGA-II \cite{deb2002fast}, as is standard for coverage--cost placement in WSN and radar management \cite{molina2008layout,iqbal2015wsn,charlish2016radar}, implemented in DEAP \cite{deap2012}. The airport study also runs two single-layer profiles, (\(M_{wp,\mathrm{coop}}\), CapEx) and (\(M_{wp,\mathrm{noncoop}}\), CapEx); the latter excludes RF from the type set. Operators:

\textbf{Parent selection}: Crowded-comparison tournament (DEAP \texttt{selTournamentDCD}).

\textbf{Crossover}: Two-point crossover on the gene list, probability 0.7.

\textbf{Mutation}: Applied with probability 0.2 per offspring; each gene is perturbed with probability 0.1, and within a perturbed gene the site, type, azimuth, and activation bit change with probabilities 0.3, 0.3, 0.25, and 0.5. Genes are then reactivated if fewer than the minimum sensor count remain active.

\textbf{Survival}: NSGA-II non-dominated sorting with crowding distance (\texttt{selNSGA2}).

Objective vectors are evaluated in parallel across CPU cores.

Budget: MUSCAT's exhaustive search grows as \(O(2^n)\) in the number of candidate sites \(n\), which rules it out for the synthetic city (221 sites, four types, up to 15 sensors). A city run uses population 32 and 18 generations, at most \(32\times19=608\) evaluations. The fronts we report are the non-dominated sets found within that budget, not certified Pareto-optimal sets; Section~\ref{sec:results_seeds} compares them with uniform random search at the same budget.


